\subsection{The dyadic evaluation on the harmonic quartic cover}

\begin{statement}{A.111}{Proved}
{Every harmonic negative-triple character has zero dyadic evaluation}
Put $F(z)=z(z-1)(z-\tfrac12)$.
Let $R,T,c\in\Q_2\setminus\{0,1,\tfrac12\}$, with
$c\ne R,T$, and suppose that
\[
 \frac{T}{R},\qquad\frac{T-1}{R-1},\qquad
                  \frac{T-\tfrac12}{R-\tfrac12}
                         \quad\hbox{belong to }\Q_2^{\times4}.
\tag{A.111.1}
\]
For every $r\in\Q_2^\times$ satisfying
\[
                 r^2=-\frac{F(c)}{F(R)},
\tag{A.111.2}
\]
the actual quaternion evaluation is
\[
       \left(\frac{T-c}{R-c},\,2r\right)=0
                                      \quad\hbox{in }\operatorname{Br}(\Q_2).
\tag{A.111.3}
\]
In particular this applies to every ordinary dyadic point
on the harmonic quartic cover of A.108 and every rational
negative-triple fibre character.  The assertion concerns
the displayed evaluation; no vanishing of a global
Jacobian cohomology class is inferred from it.
\end{statement}
\begin{proof}
Write $v=v_2$ and $f=(T-c)/(R-c)$.
We first prove the coordinate reduction by its exact maps.
The two transformations
\[
             g_1(z)=1-z,\qquad g_2(z)=\frac1{2z}
\tag{A.111.4}
\]
permute the four punctures $0,1,\tfrac12,\infty$.
Under $g_1$ the first two ratios in (A.111.1) are
exchanged and the third is unchanged.
Under $g_2$ their respective new values are
\[
 \frac RT,\qquad
 \frac RT\frac{T-\tfrac12}{R-\tfrac12},\qquad
                         \frac RT\frac{T-1}{R-1}.
\tag{A.111.5}
\]
Thus these are actual changes of the fourth-root covers:
the root coordinates are sent to the indicated inverses
and quotients, and all three fourth-power conditions
are preserved.

The polynomial and quaternion entries satisfy
\[
 \begin{aligned}
  F(g_1z)&=-F(z),&
  \frac{g_1T-g_1c}{g_1R-g_1c}&=f,\\
  F(g_2z)&=\frac{F(z)}{4z^4},&
  \frac{g_2T-g_2c}{g_2R-g_2c}&=f\frac RT.
 \end{aligned}
\tag{A.111.6}
\]
For $g_1$ use the same $r$.
For $g_2$ use $r'=r(R/c)^2$ in (A.111.2).
The second quaternion entry is then multiplied by a
square and the first by the fourth power $R/T$.
Consequently (A.111.3) is unchanged as an equality
of Brauer classes, not merely as a valuation condition.

These maps put $R$ in the disc $v(R)=e\ge1$.
If $v(R)=0$, use $g_1$; if $v(R)\le-2$, use $g_2$;
and if $v(R)=-1$, use $g_1g_2$.
All the evaluations remain defined because the
transformations permute precisely the four excluded
punctures and preserve distinctness of $R,T,c$.

\emph{The fourth-power conditions keep $T$ in this disc.}
From (A.111.1),
\[
 v(T)\equiv e,\qquad
 v(T-1)\equiv0,\qquad
 v(T-\tfrac12)\equiv-1\pmod4.
\tag{A.111.7}
\]
If $v(T)<-1$, the latter two valuations are both
$v(T)$, contradicting (A.111.7).
If $v(T)=-1$, the valuation of $T-1$ is $-1$,
again a contradiction.
If $v(T)=0$, then $v(T-1)$ is a positive multiple
of four and $e$ is a positive multiple of four.
It follows that
\[
       \frac{T-\tfrac12}{R-\tfrac12}
               =\frac{1-2T}{1-2R}\equiv-1\pmod{16},
\]
which cannot be a fourth power of a dyadic unit.
Thus $e_T=v(T)>0$.

A dyadic unit is a fourth power exactly when it
is one modulo sixteen.
The second ratio of (A.111.1) therefore gives
$v(T-R)\ge4$.
The first ratio gives $e_T\equiv e\pmod4$ and,
when $e_T=e$, the stronger inequality
$v(T-R)\ge e+4$.
In particular
\[
 \begin{array}{ll}
 1\le e\le3:& e_T=e,\quad v(T-R)\ge e+4,\\
 e\ge4:&e_T\ge4,\quad e_T\equiv e\pmod4.
 \end{array}
\tag{A.111.8}
\]
The first line follows because unequal positive
valuations would give $v(T-R)=\min(e,e_T)<4$;
the same argument excludes $e_T<4$ in the second line.

Write
\[
       F(z)=\tfrac12 zG(z),\qquad
                      G(z)=(1-z)(1-2z)=1-3z+2z^2.
\tag{A.111.9}
\]
For $v(z)>0$, $G(z)$ is a unit.
If $v(c)\le0$, then $v(R-c)=v(c)\le0$,
and (A.111.8) gives
$f=1+(T-R)/(R-c)\in1+16\Z_2$.
It is a square, which proves (A.111.3) in this case.
We henceforth write $k=v(c)>0$.

\emph{The shallow discs.}
Suppose $1\le e\le3$.
If $k\ne e$, then $v(R-c)=\min(e,k)\le e$;
thus $f\in1+16\Z_2$ by (A.111.8).
If $k=e$, put $c=Rz$ with $z$ a unit.
The alternative $v(z-1)\ge2$ is incompatible
with (A.111.2).
Indeed
\[
 G(Rz)-G(R)=R(z-1)\bigl(-3+2R(z+1)\bigr)
\tag{A.111.10}
\]
has valuation at least $e+2\ge3$.
Consequently
$r^2=-zG(Rz)/G(R)\equiv-z\pmod8$.
But $z\equiv1\pmod4$ would make this residue
either three or seven, which is not a unit square.
Therefore $v(z-1)\le1$ and
$v(R-c)\le e+1$.
Equation (A.111.8) now gives $f\in1+8\Z_2$.
Every such unit is a dyadic square, proving the
assertion throughout these three shallow discs.

\emph{Two deep coordinates and an explicit norm.}
Suppose $e\ge4$ and $k\ge4$.
Then $G(c)/G(R)\in1+16\Z_2$.
It has a square root $l\in1+8\Z_2$; this $l$
is itself a square in $\Q_2$.
Set $a=r/l$.  Equations (A.111.2) and (A.111.9)
give $a^2=-c/R$, and $2a$ has the same squareclass
as $2r$.
Choose the actual root $X\in\Q_2^\times$ with
$T/R=X^4$, supplied by (A.111.1).
Then
\[
 \begin{split}
 f&=\frac{X^4+a^2}{1+a^2},\\
 X^4+a^2&=(X^2+a)^2-(2a)X^2,\\
 1+a^2&=(1+a)^2-2a.
 \end{split}
\tag{A.111.11}
\]
The last two expressions are literal norms from
the quadratic algebra $\Q_2(\sqrt{2a})$.
Their nonzero ratio is a norm as well.
This proves (A.111.3), including the split
quadratic-algebra case.

\emph{The remaining valuations.}
Finally take $e\ge4$ and $1\le k\le3$.
From (A.111.2) and the unit factors in (A.111.9),
\[
           v(r)=\frac{k-e}{2},\qquad k\equiv e\pmod2.
\tag{A.111.12}
\]
In particular
$\min(e,e_T)-k\ge2$: for $k=2$ both larger
valuations are even and at least four; for $k=1,3$
they are odd and at least five.
Since $v(R-c)=k$, this proves $f\in1+4\Z_2$.
If $v(f-1)\ge3$, $f$ is a square and there is
nothing more to prove.
Otherwise $v(T-R)=k+2$.
The equality $e=e_T$ is impossible by the stronger
bound in (A.111.8), so
\[
 \min(e,e_T)=k+2,
 \qquad e\equiv k+2\pmod4,
 \qquad v(2r)=1+\frac{k-e}{2}\equiv0\pmod2.
\tag{A.111.13}
\]
In this last case $f\equiv5\pmod8$, so $f/5$ is
a square and the actual quadratic extension is
$L=\Q_2(\sqrt f)=\Q_2(\sqrt5)$.
It is unramified: an integral generator satisfies
$X^2-X-1$, whose reduction $X^2+X+1$ is
irreducible and separable over $\mathbf F_2$.
Its exact norm group is
\[
       N_{L/\Q_2}(L^\times)
              =\{b\in\Q_2^\times:v_2(b)\in2\Z\}.
\tag{A.111.14}
\]
Indeed a uniformizer has norm of valuation two.
The residue-field norm is onto, and on successive
principal-unit layers the norm is the surjective
trace $\mathbf F_4\to\mathbf F_2$.
Successive correction in these complete unit groups
therefore gives every unit as an actual norm.
Equation (A.111.13) places $2r$ in (A.111.14),
so the quaternion $(f,2r)$ is split.
This completes every case in (A.111.3).
\end{proof}

The exact certificate
\path{code/harmonic_dyadic_norm_certificate.py},
with its retained output, checks the displayed
coordinate identities and norm polynomial identities,
and evaluates finite families of actual dyadic
fourth-power fibres and locally soluble negative-triple
sources.  The valuation argument above supplies the
unbounded statement.
